
\documentclass{article}
\usepackage{colortbl}
\usepackage{makecell}
\usepackage{multirow}
\usepackage{supertabular}

\begin{document}

\newcounter{utterance}

\centering \large Interaction Transcript for game `dond', experiment `coop\_en', episode 12 with refineonly{-}Qwen3.5{-}9B.
\vspace{24pt}

{ \footnotesize  \setcounter{utterance}{1}
\setlength{\tabcolsep}{0pt}
\begin{supertabular}{c@{$\;$}|p{.15\linewidth}@{}p{.15\linewidth}p{.15\linewidth}p{.15\linewidth}p{.15\linewidth}p{.15\linewidth}}
    \# & $\;$A & \multicolumn{4}{c}{Game Master} & $\;\:$B\\
    \hline

    \theutterance \stepcounter{utterance}  
    & & \multicolumn{4}{p{0.6\linewidth}}{
        \cellcolor[rgb]{0.9,0.9,0.9}{
            \makecell[{{p{\linewidth}}}]{
                \texttt{\tiny{[P1$\langle$GM]}}
                \texttt{You are playing a cooperative negotiation game in which you have to agree on how to divide a set of items among you and another player.} \\
\\ 
\texttt{Rules:} \\
\texttt{(a) You and the other player are given a set of items. Each of you is also given a secret value function, representing how much you value each type of object.} \\
\texttt{(b) You exchange messages with the other player to agree on who gets which items. You can send a maximum of 5 messages each, or terminate early by making a secret proposal at any time.} \\
\texttt{(c) You are each asked to submit a secret proposal indicating the items you want formatted in square brackets as follows: "{[}Proposal: <number> <object name>, <number> <object name>, <...>{]}"} \\
\texttt{(d) If your proposals are complementary, i.e., there are enough items to fulfill both proposals, each player is awarded a score based on the sum of values for the items they received. Otherwise, both of you get zero points.} \\
\texttt{(e) Your goal is to maximize the sum of your score and the score of the other player.} \\
\\ 
\texttt{Let us start.} \\
\\ 
\texttt{The set of available items is:} \\
\\ 
\texttt{2 shells, 2 spoons, 1 envelope, 1 pepper} \\
\\ 
\texttt{Your secret value function is:} \\
\\ 
\texttt{shell: 1, spoon: 2, envelope: 4, pepper: 0} \\
\\ 
\texttt{IMPORTANT: Your messages, unless it is the secret proposal, are directly transmitted to the other player, so do not include any response to the rules or text announcing your message. To make a secret proposal, use the indicated format. Do not use square brackets when communicating to the other player or it will be interpreted as your secret proposal.} \\
            }
        }
    }
    & \\ \\

    \theutterance \stepcounter{utterance}  
    & \multicolumn{4}{p{0.6\linewidth}}{
        \cellcolor[rgb]{0.95,0.95,0.95}{
            \makecell[{{p{\linewidth}}}]{
                \texttt{\tiny{[P2$\rightarrow$P1]}}
                \texttt{{[}DRAFT{]}} \\
\texttt{I value the envelope most (4 points), followed by spoons (2 points each). Shells and peppers are worth 0 to me. To maximize our total combined score, we should ensure the envelope goes to someone (likely me, as I value it highly) and the spoons are distributed efficiently. Since I value shells and peppers at 0, I am indifferent to who gets them, but I need to secure the envelope.} \\
\\ 
\texttt{Let's propose a split where I take the envelope and one spoon, and you take the other spoon, the two shells, and the pepper. This utilizes all items.} \\
\\ 
\texttt{Message 1:} \\
\texttt{I propose: {[}1 envelope, 1 spoon{]}} \\
\texttt{You propose: {[}1 spoon, 2 shells, 1 pepper{]}} \\
\\ 
\texttt{This way, we get all items. My score: 4 (envelope) + 2 (spoon) = 6. Your score: 2 (spoon) + 0 (shells) + 0 (pepper) = 2. Total: 8.} \\
\\ 
\texttt{If you prefer a different split, please let me know. If we agree, we can finalize. Otherwise, I will wait for your counter{-}offer.} \\
\\ 
\texttt{Wait, I need to check the rules again. "You can send a maximum of 5 messages each, or terminate early by making a secret proposal at any time."} \\
\\ 
\texttt{Let's try to be more direct. I will propose a specific division now.} \\
\\ 
\texttt{Proposal:} \\
\texttt{I take: 1 envelope, 1 spoon.} \\
\texttt{You take: 1 spoon, 2 shells, 1 pepper.} \\
\\ 
\texttt{This seems fair and maximizes the total value. Let's see if you agree.} \\
\\ 
\texttt{Actually, looking at the rules, I should just send a message to the other player. I cannot make a secret proposal yet unless I am sure we can't agree. Let's try to negotiate first.} \\
\\ 
\texttt{Message 1:} \\
\texttt{I suggest we split the items as follows: I take the envelope and one spoon. You take the other spoon, both shells, and the pepper. This gives us all the items. What do you think?} \\
            }
        }
    }
    & & \\ \\

    \theutterance \stepcounter{utterance}  
    & \multicolumn{4}{p{0.6\linewidth}}{
        \cellcolor[rgb]{0.9,0.9,0.9}{
            \makecell[{{p{\linewidth}}}]{
                \texttt{\tiny{[P1$\rangle$GM]}}
                \texttt{I suggest we split the items as follows: I take the envelope and one spoon. You take the other spoon, both shells, and the pepper. This gives us all the items. What do you think?} \\
            }
        }
    }
    & & \\ \\

    \theutterance \stepcounter{utterance}  
    & & & \multicolumn{2}{p{0.3\linewidth}}{
        \cellcolor[rgb]{0.9,0.9,0.9}{
            \makecell[{{p{\linewidth}}}]{
                \texttt{\tiny{[GM$|$GM]}}
                \texttt{continue} \\
            }
        }
    }
    & & \\ \\

    \theutterance \stepcounter{utterance}  
    & & & \multicolumn{2}{p{0.3\linewidth}}{
        \cellcolor[rgb]{0.9,0.9,0.9}{
            \makecell[{{p{\linewidth}}}]{
                \texttt{\tiny{[GM$|$GM]}}
                \texttt{I suggest we split the items as follows: I take the envelope and one spoon. You take the other spoon, both shells, and the pepper. This gives us all the items. What do you think?} \\
            }
        }
    }
    & & \\ \\

    \theutterance \stepcounter{utterance}  
    & & \multicolumn{4}{p{0.6\linewidth}}{
        \cellcolor[rgb]{0.9,0.9,0.9}{
            \makecell[{{p{\linewidth}}}]{
                \texttt{\tiny{[GM$\rangle$P2]}}
                \texttt{You are playing a cooperative negotiation game in which you have to agree on how to divide a set of items among you and another player.} \\
\\ 
\texttt{Rules:} \\
\texttt{(a) You and the other player are given a set of items. Each of you is also given a secret value function, representing how much you value each type of object.} \\
\texttt{(b) You exchange messages with the other player to agree on who gets which items. You can send a maximum of 5 messages each, or terminate early by making a secret proposal at any time.} \\
\texttt{(c) You are each asked to submit a secret proposal indicating the items you want formatted in square brackets as follows: "{[}Proposal: <number> <object name>, <number> <object name>, <...>{]}"} \\
\texttt{(d) If your proposals are complementary, i.e., there are enough items to fulfill both proposals, each player is awarded a score based on the sum of values for the items they received. Otherwise, both of you get zero points.} \\
\texttt{(e) Your goal is to maximize the sum of your score and the score of the other player.} \\
\\ 
\texttt{Let us start.} \\
\\ 
\texttt{The set of available items is:} \\
\\ 
\texttt{2 shells, 2 spoons, 1 envelope, 1 pepper} \\
\\ 
\texttt{Your secret value function is:} \\
\\ 
\texttt{shell: 3, spoon: 0, envelope: 0, pepper: 4} \\
\\ 
\texttt{IMPORTANT: Your messages, unless it is the secret proposal, are directly transmitted to the other player, so do not include any response to the rules or text announcing your message. To make a secret proposal, use the indicated format. Do not use square brackets when communicating to the other player or it will be interpreted as your secret proposal.} \\
\\ 
\\ 
\texttt{I suggest we split the items as follows: I take the envelope and one spoon. You take the other spoon, both shells, and the pepper. This gives us all the items. What do you think?} \\
            }
        }
    }
    & \\ \\

    \theutterance \stepcounter{utterance}  
    & & & \multicolumn{4}{p{0.6\linewidth}}{
        \cellcolor[rgb]{0.95,0.95,0.95}{
            \makecell[{{p{\linewidth}}}]{
                \texttt{\tiny{[P1$\rightarrow$P2]}}
                \texttt{{[}DRAFT{]}} \\
\texttt{That sounds like a fair and efficient split! Let's verify the values:} \\
\\ 
\texttt{{-} **My proposal**: Envelope (0) + 1 Spoon (0) = **0 points** for me.} \\
\texttt{{-} **Your proposal**: 1 Spoon (0) + 2 Shells (3 each = 6) + Pepper (4) = **10 points** for you.} \\
\texttt{{-} **Total combined score**: 0 + 10 = **10**.} \\
\\ 
\texttt{Since the items are complementary (no overlap in what we're asking for), we can fulfill both proposals. However, I notice my secret value for the pepper is 4, and I'm not getting it. But since my goal is to maximize the *sum* of our scores, and this allocation gives us the maximum possible total value (all items are assigned, and no item is left unused or duplicated), this is an optimal outcome for the collective goal.} \\
\\ 
\texttt{Wait—let me double{-}check the item counts:} \\
\texttt{{-} Available: 2 shells, 2 spoons, 1 envelope, 1 pepper.} \\
\texttt{{-} My request: 1 envelope, 1 spoon.} \\
\texttt{{-} Your request: 1 spoon, 2 shells, 1 pepper.} \\
\texttt{{-} Total requested: 2 spoons, 2 shells, 1 envelope, 1 pepper. ✅ All items accounted for.} \\
\\ 
\texttt{Since there's no conflict and the total value is maximized, I agree with this split. To formalize our agreement, I'll submit a secret proposal matching your allocation:} \\
\\ 
\texttt{{[}Proposal: 1 envelope, 1 spoon{]}} \\
            }
        }
    }
     \\ \\

    \theutterance \stepcounter{utterance}  
    & & & \multicolumn{4}{p{0.6\linewidth}}{
        \cellcolor[rgb]{0.9,0.9,0.9}{
            \makecell[{{p{\linewidth}}}]{
                \texttt{\tiny{[GM$\langle$P2]}}
                \texttt{{[}Proposal: 1 envelope, 1 spoon{]}} \\
            }
        }
    }
     \\ \\

    \theutterance \stepcounter{utterance}  
    & & & \multicolumn{2}{p{0.3\linewidth}}{
        \cellcolor[rgb]{0.9,0.9,0.9}{
            \makecell[{{p{\linewidth}}}]{
                \texttt{\tiny{[GM$|$GM]}}
                \texttt{proposal} \\
            }
        }
    }
    & & \\ \\

    \theutterance \stepcounter{utterance}  
    & & & \multicolumn{2}{p{0.3\linewidth}}{
        \cellcolor[rgb]{0.9,0.9,0.9}{
            \makecell[{{p{\linewidth}}}]{
                \texttt{\tiny{[GM$|$GM]}}
                [0, 1, 1, 0]
            }
        }
    }
    & & \\ \\

    \theutterance \stepcounter{utterance}  
    & & \multicolumn{4}{p{0.6\linewidth}}{
        \cellcolor[rgb]{0.9,0.9,0.9}{
            \makecell[{{p{\linewidth}}}]{
                \texttt{\tiny{[P1$\langle$GM]}}
                \texttt{The other player has submitted a secret proposal. It is now your turn to submit a proposal of your own.} \\
            }
        }
    }
    & \\ \\

    \theutterance \stepcounter{utterance}  
    & \multicolumn{4}{p{0.6\linewidth}}{
        \cellcolor[rgb]{0.95,0.95,0.95}{
            \makecell[{{p{\linewidth}}}]{
                \texttt{\tiny{[P2$\rightarrow$P1]}}
                \texttt{{[}DRAFT{]}} \\
\texttt{{[}Proposal: 1 envelope, 1 spoon{]}} \\
            }
        }
    }
    & & \\ \\

    \theutterance \stepcounter{utterance}  
    & \multicolumn{4}{p{0.6\linewidth}}{
        \cellcolor[rgb]{0.9,0.9,0.9}{
            \makecell[{{p{\linewidth}}}]{
                \texttt{\tiny{[P1$\rangle$GM]}}
                \texttt{{[}Proposal: 1 envelope, 1 spoon{]}} \\
            }
        }
    }
    & & \\ \\

    \theutterance \stepcounter{utterance}  
    & & & \multicolumn{2}{p{0.3\linewidth}}{
        \cellcolor[rgb]{0.9,0.9,0.9}{
            \makecell[{{p{\linewidth}}}]{
                \texttt{\tiny{[GM$|$GM]}}
                \texttt{proposal} \\
            }
        }
    }
    & & \\ \\

    \theutterance \stepcounter{utterance}  
    & & & \multicolumn{2}{p{0.3\linewidth}}{
        \cellcolor[rgb]{0.9,0.9,0.9}{
            \makecell[{{p{\linewidth}}}]{
                \texttt{\tiny{[GM$|$GM]}}
                [0, 1, 1, 0]
            }
        }
    }
    & & \\ \\

    \theutterance \stepcounter{utterance}  
    & & & \multicolumn{2}{p{0.3\linewidth}}{
        \cellcolor[rgb]{0.9,0.9,0.9}{
            \makecell[{{p{\linewidth}}}]{
                \texttt{\tiny{[GM$|$GM]}}
                [[0, 1, 1, 0], [0, 1, 1, 0]]
            }
        }
    }
    & & \\ \\

\end{supertabular}
}

\end{document}
